Thus:
@@M0@@
Now:
@@M1@@
Putting @@M2@@, we therefore obtain:
@@M3@@
Moreover, @@M4@@; now @@M5@@; thus @@M6@@.

\sgaStableTarget{sga6.e06.prop.5-11}\textbf{Proposition 5.11.} \uline{With the notation and hypotheses of} 5.1, \uline{let} @@M7@@. \uline{For} @@M8@@ \uline{to hold, it is necessary and sufficient that there exist} @@M9@@ \uline{such that}:
@@M10@@
The proof is identical to that of 5.10, working in the associated graded rings and using 5.6.

\sgaStableTarget{sga6.e06.sec.6}\section*{6. Study of the filtration of \texorpdfstring{@@M11@@}{K^{\bullet}(X)}, @@M12@@ having an ample sheaf}

In this section, @@M13@@ is a scheme having an ample sheaf @@M14@@. Recall that in this case the @@M15@@ defined in terms of locally free sheaves of finite type, or of perfect complexes, are equal (IV, 2.9). We denote by
